% concatenated from A_nonlocal_transmission_problem_on_a_hybrid_continuous-discrete_domain.tex
\documentclass[11pt]{article}

\usepackage[T1]{fontenc}
\usepackage[utf8]{inputenc}
\usepackage{lmodern}
\usepackage{amsmath,amssymb,amsthm,mathtools}
\usepackage[a4paper,margin=1in]{geometry}
\usepackage{xcolor}
\usepackage{hyperref}
\usepackage{mathrsfs}
\usepackage{enumitem}
\numberwithin{equation}{section}
\hypersetup{colorlinks=true, linkcolor=blue!50!black, citecolor=blue!50!black, urlcolor=blue!50!black}

\newtheorem{theorem}{Theorem}[section]
\newtheorem{proposition}[theorem]{Proposition}
\newtheorem{lemma}[theorem]{Lemma}
\newtheorem{corollary}[theorem]{Corollary}
\theoremstyle{definition}
\newtheorem{definition}[theorem]{Definition}
\newtheorem{remark}[theorem]{Remark}

\newcommand{\R}{\mathbb{R}}
\newcommand{\T}{\mathbb{T}}
\newcommand{\EE}{\mathcal{E}}
\newcommand{\JJ}{\mathcal{J}}
\newcommand{\HH}{\mathcal{H}}
\newcommand{\ddmu}{\,d\mu}
\newcommand{\abs}[1]{\lvert#1\rvert}
\newcommand{\norm}[1]{\lVert#1\rVert}
\newcommand{\PV}{\operatorname{P.V.}}
\newcommand{\cc}{\mathrm{cc}}
\newcommand{\cdt}{\mathrm{cd}}
\newcommand{\dd}{\mathrm{dd}}
\newcommand{\subjclass}[2][]{\par\smallskip\noindent\textbf{MSC 2020: }#2\par}
\newcommand{\keywords}[1]{\par\smallskip\noindent\textbf{Keywords: }#1\par}

\title{\textbf{A nonlocal transmission problem on a hybrid continuous-discrete domain}}
\author{Hafida Abbas\thanks{Department of Economics, Faculty of Economics and Commercial Sciences, University of Saida, Algeria.} \and Abdelhalim Azzouz\thanks{Department of Mathematics, Faculty of Mathematics, Computer Science and Telecoms, University of Saida, Algeria. Corresponding author: \texttt{abdelhalim.azzouz@univ-saida.dz}.}}
\date{}

\begin{document}
\maketitle

\begin{abstract}
We study a quadratic nonlocal variational problem on a hybrid domain formed by a compact interval and finitely many discrete points. The associated energy splits into continuous, discrete, and interface contributions. Our main estimate shows that the interface term yields a coercive coupling between the two phases and provides an equivalent hybrid norm. As a consequence, we prove existence and uniqueness of a minimizer for the corresponding variational problem and characterize it as the unique weak solution of the associated hybrid Euler--Lagrange system. The latter combines a nonlocal integral equation on the continuous component with a finite nonlocal algebraic system on the discrete nodes.
\end{abstract}

\subjclass[2020]{Primary 26A33, 35R11, 46E35; Secondary 34B37, 34N05, 49J45.}
\keywords{nonlocal variational problem, hybrid continuous-discrete domain, interface coercivity, weak solution, Euler--Lagrange system.}

\section{Introduction}

Fractional energies arise naturally in the analysis of nonlocal equations, variational problems, and regularity theory. In the standard continuous setting they are generated by the Gagliardo seminorm and lead to the classical fractional Sobolev spaces; see, for instance, \cite{DiNezza2012, Gagliardo1957, Slobodeckij1958}. On the other hand, the calculus on time scales was introduced to provide a common language for continuous and discrete models; see \cite{Hilger1990, BohnerPeterson2001, BohnerPeterson2003}.

From the measure-theoretic point of view, the time-scale framework already provides a natural notion of Lebesgue $\Delta$-measure and multidimensional $\Delta$-integration. In particular, Bohner and Guseinov developed a multiple Lebesgue integration theory on products of time scales, while Cabada and Vivero showed how the Lebesgue $\Delta$-integral can be expressed in terms of an ordinary Lebesgue integral on the real line; see \cite{BohnerGuseinov2006, CabadaVivero2006}. These results justify the use of measure-based function spaces on $\T$ and on $\T\times\T$.

On this basis, first-order Lebesgue and Sobolev spaces on time scales have already been studied in several directions. Basic properties of Sobolev spaces endowed with the Lebesgue $\Delta$-measure were established in \cite{AgarwalOteroPereraVivero2006}, and generalized Lebesgue and Sobolev spaces with variable exponent on bounded time scales were investigated in \cite{SkrzypekSzymanska2019}. Thus, at order one, the functional-analytic setting is already available in the literature.

The fractional situation is different. Existing fractional theories on time scales are developed predominantly through specific notions of fractional differentiation, such as conformable derivatives or Riemann--Liouville-type derivatives; see, for example, \cite{BenkhettouHassaniTorres2016, HuLi2022Right, HuLi2022Left}. By contrast, a genuinely nonlocal Gagliardo-type approach based on a double-integral energy over $\T\times\T$ appears to be much less explored on time scales, especially when one has in mind variable-order constructions. This is the point of view from which the present work originates.

The purpose of the present paper is to study a nonlocal variational problem on a simple hybrid continuous-discrete domain and to identify the coercive role of the interface contribution. More precisely, for an integer $N\ge 2$, we consider
\[
\T=[0,1]\cup\{2,3,\dots,N+1\},
\]
endowed with the measure
\[
\mu=\mathbf{1}_{[0,1]}(x)\,dx+\sum_{k=2}^{N+1}\delta_k.
\]
This model consists of a continuous phase and a finite discrete phase separated by a positive gap. For functions $u:\T\to\R$, the nonlocal quadratic energy generated by the kernel $\abs{x-y}^{-1-2\alpha}$ splits into three terms:
\begin{itemize}
\item a continuous--continuous interaction on $[0,1]\times[0,1]$,
\item a discrete--discrete interaction on $\{2,\dots,N+1\}^2$,
\item a continuous--discrete interface contribution.
\end{itemize}
The first two terms correspond to the pure continuous and pure discrete components. The third one is the key feature of the model: it couples the continuous profile to the discrete variables and yields a coercive transmission mechanism between the two phases.

Our main estimate shows that the interface contribution is not a lower-order perturbation. It controls simultaneously the $L^2$ oscillation of the continuous component around its mean and the distance of each discrete value to that same mean. Consequently, the full energy is equivalent to a natural hybrid norm. This provides a robust variational framework for the boundary-value-type problem associated with the quadratic functional
\[
\JJ(u):=\frac12\EE_\alpha(u)+\frac{\lambda}{2}\norm{u}_{L^2(\T,\mu)}^2-\int_{\T}fu\ddmu,
\qquad \lambda>0.
\]
Within this framework we prove existence and uniqueness of a minimizer, define the corresponding notion of weak solution, and derive the associated hybrid Euler--Lagrange system. We interpret \eqref{eq:weak-solution} as the weak formulation of a hybrid nonlocal transmission problem on the continuous-discrete domain $\T$. The latter consists of a nonlocal integral equation on the continuous phase coupled with a finite nonlocal algebraic system on the discrete nodes.

The paper is organized as follows. Section~\ref{sec:coercive-framework} introduces the hybrid measure space, the energy, its decomposition, and the main coercivity estimate. Section~\ref{sec:variational} formulates the variational problem and establishes the equivalence between minimizers and weak solutions, together with existence and uniqueness. Section~\ref{sec:transmission} derives the hybrid transmission system in weak form and discusses its formal pointwise counterpart. The final section presents a concrete variational example on a hybrid time scale and illustrates the continuous-discrete transmission structure of the model.

\section{The hybrid setting and the main coercivity result}\label{sec:coercive-framework}
While the measure-theoretic and first-order Sobolev frameworks on time scales are already available, our purpose here is different: we introduce a nonlocal Gagliardo-type energy on the hybrid continuous-discrete domain under consideration and use it to formulate the transmission problem studied in the sequel.
\subsection{The hybrid domain and the nonlocal energy}

Fix an integer $N\ge 2$ and define
\[
\T=[0,1]\cup D,
\qquad D:=\{2,3,\dots,N+1\}.
\]
We equip $\T$ with the positive finite Borel measure
\begin{equation}\label{eq:def-mu}
\mu(A):=\int_{A\cap[0,1]}dx+\sum_{k=2}^{N+1}\delta_k(A),
\end{equation}
for every Borel set $A\subset \T$. Equivalently,
\[
\int_{\T} f\ddmu =\int_0^1 f(x)\,dx+\sum_{k=2}^{N+1} f(k)
\]
for every integrable function $f$ on $\T$.

A function $u:\T\to\R$ will be written as
\[
u=(v,a_2,\dots,a_{N+1}),
\]
where
\[
v=u|_{[0,1]},\qquad a_k=u(k)\quad (k=2,\dots,N+1).
\]
Thus the continuous and discrete phases are encoded separately.

Let $\alpha\in(0,1)$. For $u\in L^2(\T,\mu)$, define the quadratic energy
\begin{equation}\label{eq:def-energy}
\EE_\alpha(u):=
\iint_{\T\times\T}
\frac{\abs{u(x)-u(y)}^2}{\abs{x-y}^{1+2\alpha}}
\,d\mu(x)\,d\mu(y).
\end{equation}
Since the discrete phase is finite and separated from $[0,1]$ by a positive distance, the mixed and discrete contributions are automatically finite whenever the restriction $v$ belongs to $H^\alpha(0,1)$.

\begin{definition}
We define the hybrid energy space by
\[
\HH_\alpha(\T):=
\Big\{u=(v,a_2,\dots,a_{N+1}): v\in H^\alpha(0,1),\ a_k\in\R\Big\}.
\]
It is endowed with the norm
\begin{equation}\label{eq:def-Hnorm}
\norm{u}_{\HH_\alpha(\T)}^2:=\norm{u}_{L^2(\T,\mu)}^2+\EE_\alpha(u).
\end{equation}
\end{definition}

The continuous--continuous part of \eqref{eq:def-energy} is the standard Gagliardo seminorm squared on $(0,1)$, namely
\begin{equation}\label{eq:gagliardo-seminorm}
[v]_{H^\alpha(0,1)}^2:=\int_0^1\int_0^1\frac{\abs{v(x)-v(y)}^2}{\abs{x-y}^{1+2\alpha}}\,dx\,dy.
\end{equation}
\begin{lemma}[Equivalence with the product norm]\label{lem:product-norm}
There exist constants $c_0,C_0>0$, depending only on $\alpha$ and $N$, such that for every
\[u=(v,a_2,\dots,a_{N+1})\in \HH_\alpha(\T),
\]
one has
\[
c_0\Bigl(\norm{v}_{H^\alpha(0,1)}^2+\sum_{k=2}^{N+1}\abs{a_k}^2\Bigr)
\le \norm{u}_{\HH_\alpha(\T)}^2
\le C_0\Bigl(\norm{v}_{H^\alpha(0,1)}^2+\sum_{k=2}^{N+1}\abs{a_k}^2\Bigr).
\]
Consequently, $\HH_\alpha(\T)$ is a Hilbert space.
\end{lemma}

\begin{proof}
The upper bound follows from the decomposition in Proposition~\ref{prop:decomposition}. The continuous--continuous term is exactly $[v]_{H^\alpha(0,1)}^2$. Since the discrete phase is finite and satisfies $\operatorname{dist}([0,1],D)=1$, one has
\[
1\le \abs{k-x}\le N+1\qquad (x\in[0,1],\ k\in D),
\]
which gives
\[
\EE_{\cdt}(v,a)
\le \sum_{k=2}^{N+1}\int_0^1 \abs{a_k-v(x)}^2\,dx
\le 2\sum_{k=2}^{N+1}\abs{a_k}^2+2N\norm{v}_{L^2(0,1)}^2.
\]
Similarly, since $D$ is finite,
\[
\EE_{\dd}(a)
\le 2\sum_{\substack{i,j=2\\ i\neq j}}^{N+1}(\abs{a_i}^2+\abs{a_j}^2)
\le 4N\sum_{k=2}^{N+1}\abs{a_k}^2.
\]
Hence
\[
\EE_\alpha(u)\le C\Bigl([v]_{H^\alpha(0,1)}^2+\norm{v}_{L^2(0,1)}^2+\sum_{k=2}^{N+1}\abs{a_k}^2\Bigr)
\]
for some $C>0$, which yields the upper bound.

For the lower bound, we simply note that
\[
\norm{u}_{\HH_\alpha(\T)}^2=\norm{u}_{L^2(\T,\mu)}^2+\EE_\alpha(u)
\ge \norm{u}_{L^2(\T,\mu)}^2
=\norm{v}_{L^2(0,1)}^2+\sum_{k=2}^{N+1}\abs{a_k}^2
\]
and also $\EE_\alpha(u)\ge [v]_{H^\alpha(0,1)}^2$. Combining these estimates gives
\[
\norm{u}_{\HH_\alpha(\T)}^2
\ge \norm{v}_{L^2(0,1)}^2+[v]_{H^\alpha(0,1)}^2+\sum_{k=2}^{N+1}\abs{a_k}^2
=\norm{v}_{H^\alpha(0,1)}^2+\sum_{k=2}^{N+1}\abs{a_k}^2,
\]
which proves the claim.
\end{proof}

\subsection{Decomposition of the energy and the interface estimate}

The energy \eqref{eq:def-energy} splits naturally into three contributions.

\begin{proposition}[Decomposition of the energy]\label{prop:decomposition}
For every $u=(v,a_2,\dots,a_{N+1})\in \HH_\alpha(\T)$,
\begin{equation}\label{eq:energy-decomposition}
\EE_\alpha(u)=\EE_{\cc}(v)+2\EE_{\cdt}(v,a)+\EE_{\dd}(a),
\end{equation}
where
\begin{align}
\EE_{\cc}(v)&:=\int_0^1\int_0^1\frac{\abs{v(x)-v(y)}^2}{\abs{x-y}^{1+2\alpha}}\,dx\,dy,
\label{eq:Ecc}\\
\EE_{\cdt}(v,a)&:=\sum_{k=2}^{N+1}\int_0^1\frac{\abs{a_k-v(x)}^2}{\abs{k-x}^{1+2\alpha}}\,dx,
\label{eq:Ecd}\\
\EE_{\dd}(a)&:=\sum_{\substack{i,j=2\\ i\neq j}}^{N+1}\frac{\abs{a_i-a_j}^2}{\abs{i-j}^{1+2\alpha}}.
\label{eq:Edd}
\end{align}
\end{proposition}

\begin{proof}
Using \eqref{eq:def-mu}, the product measure $\mu\otimes\mu$ decomposes into the sum of the continuous--continuous, continuous--discrete, discrete--continuous, and discrete--discrete parts. Substituting this decomposition into \eqref{eq:def-energy}, we obtain
\begin{align*}
\EE_\alpha(u)
&=\int_0^1\int_0^1\frac{\abs{v(x)-v(y)}^2}{\abs{x-y}^{1+2\alpha}}\,dx\,dy \\
&\quad +\sum_{k=2}^{N+1}\int_0^1\frac{\abs{a_k-v(x)}^2}{\abs{k-x}^{1+2\alpha}}\,dx
+\sum_{k=2}^{N+1}\int_0^1\frac{\abs{v(x)-a_k}^2}{\abs{x-k}^{1+2\alpha}}\,dx \\
&\quad +\sum_{\substack{i,j=2\\ i\neq j}}^{N+1}\frac{\abs{a_i-a_j}^2}{\abs{i-j}^{1+2\alpha}}.
\end{align*}
Since $\abs{k-x}=\abs{x-k}$ and $\abs{a_k-v(x)}=\abs{v(x)-a_k}$, the two mixed terms coincide, which yields \eqref{eq:energy-decomposition}.
\end{proof}

\begin{lemma}[Interface estimate]\label{lem:interface-estimate}
There exist constants $c_1,c_2>0$, depending only on $\alpha$ and $N$, such that for every $u=(v,a_2,\dots,a_{N+1})\in\HH_\alpha(\T)$,
\begin{equation}\label{eq:interface-equivalence}
c_1\sum_{k=2}^{N+1}\int_0^1\abs{a_k-v(x)}^2\,dx
\le \EE_{\cdt}(v,a)
\le c_2\sum_{k=2}^{N+1}\int_0^1\abs{a_k-v(x)}^2\,dx.
\end{equation}
\end{lemma}

\begin{proof}
Fix $k\in\{2,\dots,N+1\}$ and $x\in[0,1]$. Then
\[
k-1\le k-x\le k.
\]
Therefore,
\[
\frac{1}{k^{1+2\alpha}}\le \frac{1}{\abs{k-x}^{1+2\alpha}}\le \frac{1}{(k-1)^{1+2\alpha}}.
\]
Multiplying by $\abs{a_k-v(x)}^2$ and integrating over $(0,1)$, we obtain
\[
\frac{1}{k^{1+2\alpha}}\int_0^1 \abs{a_k-v(x)}^2\,dx
\le
\int_0^1 \frac{\abs{a_k-v(x)}^2}{\abs{k-x}^{1+2\alpha}}\,dx
\le
\frac{1}{(k-1)^{1+2\alpha}}\int_0^1 \abs{a_k-v(x)}^2\,dx.
\]
Summing with respect to $k$ yields \eqref{eq:interface-equivalence} with
\[
c_1=(N+1)^{-(1+2\alpha)},
\qquad
c_2=1.
\]
\end{proof}

\begin{lemma}[Algebraic decomposition around the mean]\label{lem:mean-decomposition}
Let
\[
\bar v:=\int_0^1 v(x)\,dx.
\]
Then, for every $k\in\{2,\dots,N+1\}$,
\begin{equation}\label{eq:single-mean-decomposition}
\int_0^1 \abs{a_k-v(x)}^2\,dx
=\abs{a_k-\bar v}^2+\norm{v-\bar v}_{L^2(0,1)}^2.
\end{equation}
Consequently,
\begin{equation}\label{eq:sum-mean-decomposition}
\sum_{k=2}^{N+1}\int_0^1 \abs{a_k-v(x)}^2\,dx
=
\sum_{k=2}^{N+1}\abs{a_k-\bar v}^2
+N\norm{v-\bar v}_{L^2(0,1)}^2.
\end{equation}
\end{lemma}

\begin{proof}
Fix $k\in\{2,\dots,N+1\}$. We write
\[
a_k-v(x)=(a_k-\bar v)+(\bar v-v(x)).
\]
Squaring and integrating over $(0,1)$, we obtain
\begin{align*}
\int_0^1\abs{a_k-v(x)}^2\,dx
&=\int_0^1 \abs{a_k-\bar v}^2\,dx
+\int_0^1\abs{v(x)-\bar v}^2\,dx \\
&\quad +2(a_k-\bar v)\int_0^1(\bar v-v(x))\,dx.
\end{align*}
Since
\[
\int_0^1(\bar v-v(x))\,dx=\bar v-\int_0^1v(x)\,dx=0,
\]
the mixed term vanishes. This proves \eqref{eq:single-mean-decomposition}. Summing over $k$ gives \eqref{eq:sum-mean-decomposition}.
\end{proof}

\subsection{Interface coercivity and the hybrid energy space}

We now derive the main coercivity statement.

\begin{theorem}[Interface coercivity]\label{thm:interface-coercivity}
There exist constants $C_1,C_2>0$, depending only on $\alpha$ and $N$, such that for every $u=(v,a_2,\dots,a_{N+1})\in\HH_\alpha(\T)$,
\begin{align}
C_1\Big(
[v]_{H^\alpha(0,1)}^2
+\norm{v-\bar v}_{L^2(0,1)}^2
+\sum_{k=2}^{N+1}\abs{a_k-\bar v}^2
\Big)
&\le \EE_\alpha(u)
\label{eq:coercivity-lower}\\
&\le
C_2\Big(
[v]_{H^\alpha(0,1)}^2
+\norm{v-\bar v}_{L^2(0,1)}^2
+\sum_{k=2}^{N+1}\abs{a_k-\bar v}^2
\Big).
\label{eq:coercivity-upper}
\end{align}
\end{theorem}

\begin{proof}
By Proposition~\ref{prop:decomposition},
\[
\EE_\alpha(u)=\EE_{\cc}(v)+2\EE_{\cdt}(v,a)+\EE_{\dd}(a).
\]
Since all three terms on the right-hand side are nonnegative, Lemmas~\ref{lem:interface-estimate} and \ref{lem:mean-decomposition} imply
\begin{align*}
\EE_\alpha(u)
&\ge \EE_{\cc}(v)+2\EE_{\cdt}(v,a) \\
&\ge [v]_{H^\alpha(0,1)}^2+2c_1\sum_{k=2}^{N+1}\int_0^1\abs{a_k-v(x)}^2\,dx \\
&= [v]_{H^\alpha(0,1)}^2
+2c_1\sum_{k=2}^{N+1}\abs{a_k-\bar v}^2
+2c_1N\norm{v-\bar v}_{L^2(0,1)}^2.
\end{align*}
This proves \eqref{eq:coercivity-lower}.

For the reverse estimate, we use again Proposition~\ref{prop:decomposition}. The term $\EE_{\cc}(v)$ is exactly $[v]_{H^\alpha(0,1)}^2$. By Lemmas~\ref{lem:interface-estimate} and \ref{lem:mean-decomposition},
\[
\EE_{\cdt}(v,a)
\le c_2\Big(
\sum_{k=2}^{N+1}\abs{a_k-\bar v}^2+N\norm{v-\bar v}_{L^2(0,1)}^2
\Big).
\]
It remains to estimate the discrete--discrete contribution. For $i\neq j$,
\[
\abs{a_i-a_j}^2
\le 2\abs{a_i-\bar v}^2+2\abs{a_j-\bar v}^2.
\]
Hence,
\begin{align*}
\EE_{\dd}(a)
&\le 2\sum_{\substack{i,j=2\\ i\neq j}}^{N+1}
\frac{\abs{a_i-\bar v}^2+\abs{a_j-\bar v}^2}{\abs{i-j}^{1+2\alpha}} \\
&\le 4\Big(\sum_{m=1}^{N-1}\frac{1}{m^{1+2\alpha}}\Big)
\sum_{k=2}^{N+1}\abs{a_k-\bar v}^2.
\end{align*}
Combining these estimates gives \eqref{eq:coercivity-upper}.
\end{proof}

\begin{remark}
Theorem~\ref{thm:interface-coercivity} shows that the interface term is not a negligible correction. It controls the $L^2$ oscillation of the continuous component and pins the discrete variables to the average value of the continuous phase. In the variational problem below, this interface term acts as a coercive transmission mechanism between the two phases.
\end{remark}

\begin{corollary}[Hybrid Poincar\'e-type inequality]\label{cor:hybrid-poincare}
There exists a constant $C_P>0$, depending only on $\alpha$ and $N$, such that for every $u=(v,a_2,\dots,a_{N+1})\in \HH_\alpha(\T)$,
\[
\norm{v-\bar v}_{L^2(0,1)}^2+\sum_{k=2}^{N+1}\abs{a_k-\bar v}^2
\le C_P\,\EE_\alpha(u).
\]
\end{corollary}

\begin{proof}
This is an immediate consequence of the lower bound \eqref{eq:coercivity-lower} in Theorem~\ref{thm:interface-coercivity}.
\end{proof}

A direct consequence is that the quadratic form
\[
\|u\|_{\ast}^2:=\|u\|_{L^2(\T,\mu)}^2+ [v]_{H^\alpha(0,1)}^2+\norm{v-\bar v}_{L^2(0,1)}^2+\sum_{k=2}^{N+1}\abs{a_k-\bar v}^2
\]
defines a norm equivalent to \eqref{eq:def-Hnorm} on $\HH_\alpha(\T)$.

\section{The variational problem and weak solutions}\label{sec:variational}
In line with the variational treatment of nonlocal problems in weak form and prove its well-posedness; see, for instance, \cite{NonlocalWeakSolutions,MRS2016,SteinbachWendland2000}. We now formulate the associated hybrid transmission problem in weak form and prove its well-posedness. The interface coercivity obtained in the previous section is the key tool in this step.
\subsection{Formulation of the problem}

Let $\lambda>0$ and let $f\in L^2(\T,\mu)$. We consider the quadratic functional
\begin{equation}\label{eq:def-functional}
\JJ(u):=\frac12\EE_\alpha(u)+\frac{\lambda}{2}\norm{u}_{L^2(\T,\mu)}^2-\int_{\T}f u\ddmu,
\qquad u\in\HH_\alpha(\T).
\end{equation}

The bilinear form naturally associated with the energy is
\begin{equation}\label{eq:def-bilinear}
\mathfrak a(u,\varphi):=
\iint_{\T\times\T}
\frac{(u(x)-u(y))(\varphi(x)-\varphi(y))}{\abs{x-y}^{1+2\alpha}}
\,d\mu(x)\,d\mu(y).
\end{equation}

\begin{definition}
We say that $u\in\HH_\alpha(\T)$ is a \emph{weak solution} of the hybrid nonlocal problem associated with $(\lambda,f)$ if
\begin{equation}\label{eq:weak-solution}
\mathfrak a(u,\varphi)+\lambda\int_{\T}u\varphi\ddmu=\int_{\T}f\varphi\ddmu
\qquad\text{for every }\varphi\in\HH_\alpha(\T).
\end{equation}
\end{definition}

Equation \eqref{eq:weak-solution} may be regarded as the weak form of a nonlocal transmission problem on the hybrid domain: the continuous phase is driven by a nonlocal integral operator, the discrete phase by a finite algebraic system, and the two parts are coupled through the interface term.

\subsection{Existence, uniqueness, and variational characterization}

\begin{theorem}[Existence and uniqueness of the minimizer]\label{thm:existence-uniqueness}
For every $\lambda>0$ and every $f\in L^2(\T,\mu)$, the functional $\JJ$ admits a unique minimizer in $\HH_\alpha(\T)$.
\end{theorem}

\begin{proof}
The mapping $u\mapsto \EE_\alpha(u)$ is quadratic and nonnegative, and the term $u\mapsto \lambda\norm{u}_{L^2(\T,\mu)}^2$ is strictly convex because $\lambda>0$. Therefore $\JJ$ is strictly convex on the Hilbert space $\HH_\alpha(\T)$.

We next prove coercivity. By Theorem~\ref{thm:interface-coercivity},
\[
\EE_\alpha(u)\ge C_1\Big(
[v]_{H^\alpha(0,1)}^2
+\norm{v-\bar v}_{L^2(0,1)}^2
+\sum_{k=2}^{N+1}\abs{a_k-\bar v}^2
\Big).
\]
It remains to control the common mean level $\bar v$. Since $\bar v=\int_0^1 v(x)\,dx$, Jensen's inequality gives
\[
\abs{\bar v}^2\le \int_0^1 \abs{v(x)}^2\,dx\le \norm{u}_{L^2(\T,\mu)}^2.
\]
Hence
\[
\norm{v}_{L^2(0,1)}^2
\le 2\norm{v-\bar v}_{L^2(0,1)}^2+2\abs{\bar v}^2
\le 2\norm{v-\bar v}_{L^2(0,1)}^2+2\norm{u}_{L^2(\T,\mu)}^2,
\]
and, for each $k\in\{2,\dots,N+1\}$,
\[
\abs{a_k}^2\le 2\abs{a_k-\bar v}^2+2\abs{\bar v}^2
\le 2\abs{a_k-\bar v}^2+2\norm{u}_{L^2(\T,\mu)}^2.
\]
Combining these inequalities with the lower bound above, we infer that there exists $c>0$ such that
\[
\EE_\alpha(u)+\lambda\norm{u}_{L^2(\T,\mu)}^2
\ge c\Bigl(\norm{v}_{H^\alpha(0,1)}^2+\sum_{k=2}^{N+1}\abs{a_k}^2\Bigr).
\]
By Lemma~\ref{lem:product-norm}, the right-hand side controls $\norm{u}_{\HH_\alpha(\T)}^2$. Therefore,
\[
\EE_\alpha(u)+\lambda\norm{u}_{L^2(\T,\mu)}^2
\ge c\norm{u}_{\HH_\alpha(\T)}^2.
\]
By Cauchy--Schwarz and Young's inequality,
\[
\left|\int_{\T}fu\ddmu\right|
\le \norm{f}_{L^2(\T,\mu)}\norm{u}_{L^2(\T,\mu)}
\le \frac{c}{4}\norm{u}_{\HH_\alpha(\T)}^2+C\norm{f}_{L^2(\T,\mu)}^2.
\]
Consequently,
\[
\JJ(u)\to +\infty\qquad\text{as }\norm{u}_{\HH_\alpha(\T)}\to\infty.
\]
Thus $\JJ$ is coercive and bounded from below.

Let $(u_n)$ be a minimizing sequence. Coercivity implies that $(u_n)$ is bounded in $\HH_\alpha(\T)$. Since $\HH_\alpha(\T)$ is a Hilbert space, there exist a subsequence, still denoted by $(u_n)$, and some $u\in\HH_\alpha(\T)$ such that $u_n\rightharpoonup u$ weakly in $\HH_\alpha(\T)$. The quadratic form $u\mapsto \EE_\alpha(u)+\lambda\norm{u}_{L^2(\T,\mu)}^2$ is weakly lower semicontinuous, whereas the linear functional $u\mapsto \int_{\T}fu\ddmu$ is weakly continuous. Therefore,
\[
\JJ(u)\le \liminf_{n\to\infty}\JJ(u_n)=\inf_{w\in\HH_\alpha(\T)}\JJ(w).
\]
Hence $u$ is a minimizer. Uniqueness follows from strict convexity.
\end{proof}

\begin{theorem}[Variational characterization of weak solutions]\label{thm:weak-var-equivalence}
A function $u\in\HH_\alpha(\T)$ is the unique minimizer of $\JJ$ if and only if it is the unique weak solution of \eqref{eq:weak-solution}.
\end{theorem}

\begin{proof}
Let $u\in\HH_\alpha(\T)$ and $\varphi\in\HH_\alpha(\T)$. For $t\in\R$,
\[
\JJ(u+t\varphi)
=\frac12\EE_\alpha(u+t\varphi)+\frac{\lambda}{2}\norm{u+t\varphi}_{L^2(\T,\mu)}^2-\int_{\T}f(u+t\varphi)\ddmu.
\]
Since the energy is quadratic,
\[
\frac{d}{dt}\Big|_{t=0}\frac12\EE_\alpha(u+t\varphi)=\mathfrak a(u,\varphi),
\]
and similarly,
\[
\frac{d}{dt}\Big|_{t=0}\frac{\lambda}{2}\norm{u+t\varphi}_{L^2(\T,\mu)}^2
=\lambda\int_{\T}u\varphi\ddmu.
\]
The derivative of the linear term is
\[
\frac{d}{dt}\Big|_{t=0}\int_{\T}f(u+t\varphi)\ddmu
=\int_{\T}f\varphi\ddmu.
\]
Hence $u$ is a critical point of $\JJ$ if and only if \eqref{eq:weak-solution} holds. Since $\JJ$ is strictly convex by Theorem~\ref{thm:existence-uniqueness}, it has at most one critical point. The conclusion follows.
\end{proof}

\section{The hybrid transmission system}\label{sec:transmission}

The weak formulation can be split into its continuous and discrete components. This yields a coupled transmission system between the two phases.

\begin{theorem}[Weak Euler--Lagrange identity]\label{thm:weak-EL}
Let $u=(v,a_2,\dots,a_{N+1})\in\HH_\alpha(\T)$ be the minimizer of $\JJ$. Then, for every $\varphi=(\psi,b_2,\dots,b_{N+1})\in\HH_\alpha(\T)$,
\begin{align}\label{eq:split-weak-form}
&\int_0^1\int_0^1
\frac{(v(x)-v(y))(\psi(x)-\psi(y))}{\abs{x-y}^{1+2\alpha}}\,dx\,dy \notag\\
&\quad +2\sum_{k=2}^{N+1}\int_0^1
\frac{(a_k-v(x))(b_k-\psi(x))}{\abs{k-x}^{1+2\alpha}}\,dx \\
&\quad +\sum_{\substack{i,j=2\\ i\neq j}}^{N+1}
\frac{(a_i-a_j)(b_i-b_j)}{\abs{i-j}^{1+2\alpha}}
+\lambda\int_{\T}u\varphi\ddmu
=\int_{\T}f\varphi\ddmu.\notag
\end{align}
\end{theorem}

\begin{proof}
Let $u=(v,a_2,\dots,a_{N+1})\in\HH_\alpha(\T)$ be the minimizer of $\JJ$. By Theorem~\ref{thm:weak-var-equivalence}, $u$ is the unique weak solution of \eqref{eq:weak-solution}. Hence, for every test function
\[
\varphi=(\psi,b_2,\dots,b_{N+1})\in\HH_\alpha(\T),
\]
one has
\begin{equation}\label{eq:proof-thm71-weak}
\mathfrak a(u,\varphi)+\lambda\int_{\T}u\varphi\,\ddmu=\int_{\T}f\varphi\,\ddmu.
\end{equation}
It therefore remains to compute the bilinear form $\mathfrak a(u,\varphi)$ in terms of the continuous and discrete components.

By definition,
\[
\mathfrak a(u,\varphi)=\iint_{\T\times\T}
\frac{(u(x)-u(y))(\varphi(x)-\varphi(y))}{\abs{x-y}^{1+2\alpha}}\,d\mu(x)\,d\mu(y).
\]
Using the decomposition of the measure
\[
d\mu(x)=\mathbf 1_{[0,1]}(x)\,dx+\sum_{k=2}^{N+1}\delta_k(x),
\]
the product measure $\mu\otimes\mu$ splits into four contributions:
continuous--continuous, continuous--discrete, discrete--continuous, and discrete--discrete. We examine them separately.

For the continuous--continuous part, since $u|_{[0,1]}=v$ and $\varphi|_{[0,1]}=\psi$, we obtain
\[
I_{cc}:=\int_0^1\int_0^1
\frac{(v(x)-v(y))(\psi(x)-\psi(y))}{\abs{x-y}^{1+2\alpha}}\,dx\,dy.
\]
For the continuous--discrete part, using $u(k)=a_k$ and $\varphi(k)=b_k$, we get
\[
I_{cd}:=\sum_{k=2}^{N+1}\int_0^1
\frac{(v(x)-a_k)(\psi(x)-b_k)}{\abs{x-k}^{1+2\alpha}}\,dx.
\]
For the discrete--continuous part, by the symmetry of the kernel and the identity $\abs{k-x}=\abs{x-k}$, we have
\[
I_{dc}:=\sum_{k=2}^{N+1}\int_0^1
\frac{(a_k-v(x))(b_k-\psi(x))}{\abs{k-x}^{1+2\alpha}}\,dx.
\]
Since
\[
(v(x)-a_k)(\psi(x)-b_k)=(a_k-v(x))(b_k-\psi(x)),
\]
it follows that $I_{cd}=I_{dc}$. Therefore the mixed contribution equals
\[
I_{cd}+I_{dc}=2\sum_{k=2}^{N+1}\int_0^1
\frac{(a_k-v(x))(b_k-\psi(x))}{\abs{k-x}^{1+2\alpha}}\,dx.
\]
Finally, for the discrete--discrete part we obtain
\[
I_{dd}:=\sum_{i=2}^{N+1}\sum_{j=2}^{N+1}
\frac{(a_i-a_j)(b_i-b_j)}{\abs{i-j}^{1+2\alpha}}.
\]
The terms with $i=j$ vanish because $(a_i-a_i)(b_i-b_i)=0$, hence
\[
I_{dd}=\sum_{\substack{i,j=2\\ i\neq j}}^{N+1}
\frac{(a_i-a_j)(b_i-b_j)}{\abs{i-j}^{1+2\alpha}}.
\]
Combining the four pieces yields
\begin{align*}
\mathfrak a(u,\varphi)
&=\int_0^1\int_0^1
\frac{(v(x)-v(y))(\psi(x)-\psi(y))}{\abs{x-y}^{1+2\alpha}}\,dx\,dy\\
&\quad+2\sum_{k=2}^{N+1}\int_0^1
\frac{(a_k-v(x))(b_k-\psi(x))}{\abs{k-x}^{1+2\alpha}}\,dx\\
&\quad+\sum_{\substack{i,j=2\\ i\neq j}}^{N+1}
\frac{(a_i-a_j)(b_i-b_j)}{\abs{i-j}^{1+2\alpha}}.
\end{align*}
Substituting this identity into \eqref{eq:proof-thm71-weak} gives exactly \eqref{eq:split-weak-form}. This completes the proof.
\end{proof}

\begin{remark}[Formal pointwise system]
Assume that the weak solution $u=(v,a_2,\dots,a_{N+1})$ has sufficient additional regularity and integrability for the pointwise expressions below to make sense. Then \eqref{eq:split-weak-form} corresponds formally to
\begin{equation}\label{eq:strong-continuous}
2\,\PV\int_0^1\frac{v(x)-v(y)}{\abs{x-y}^{1+2\alpha}}\,dy
+2\sum_{k=2}^{N+1}\frac{v(x)-a_k}{\abs{k-x}^{1+2\alpha}}
+\lambda v(x)=f(x)
\end{equation}
for $x\in(0,1)$, and to
\begin{equation}\label{eq:strong-discrete}
2\int_0^1\frac{a_k-v(x)}{\abs{k-x}^{1+2\alpha}}\,dx
+2\sum_{\substack{j=2\\ j\neq k}}^{N+1}\frac{a_k-a_j}{\abs{k-j}^{1+2\alpha}}
+\lambda a_k=f(k)
\end{equation}
for each $k\in\{2,\dots,N+1\}$.

The first equation is a nonlocal integral equation on the continuous phase, whereas the second one is a finite nonlocal algebraic system on the discrete phase. The interface contribution acts as a transmission term: it inserts the discrete variables into the continuous equation and, conversely, the continuous profile into each discrete balance law.
\end{remark}


\section{A variational example on a hybrid time scale}

In this section, we study a concrete variational problem on a hybrid time scale and derive both its weak formulation and its coupled continuous-discrete structure. Unlike the finite-dimensional reduction introduced later, this problem is formulated on the full energy space and therefore retains the full variational structure of the model.

\subsection{The hybrid time scale and the energy space}

We consider the hybrid time scale
\[
\mathbb T=[0,1]\cup\{2,3\},
\]
endowed with the measure
\[
d\mu(x)=\mathbf 1_{[0,1]}(x)\,dx+\delta_2+\delta_3.
\]
A function $u:\mathbb T\to\mathbb R$ is identified with a triple
\[
u=(v,a,b),
\]
where
\[
v=u|_{[0,1]},\qquad a=u(2),\qquad b=u(3).
\]

For $\alpha\in(0,1)$, we define the hybrid energy
\begin{align}
\mathcal E_\alpha(u)
&=
\int_0^1\int_0^1
\frac{|v(x)-v(y)|^2}{|x-y|^{1+2\alpha}}\,dx\,dy
\notag\\
&\quad
+2\int_0^1\frac{|a-v(x)|^2}{|2-x|^{1+2\alpha}}\,dx
+2\int_0^1\frac{|b-v(x)|^2}{|3-x|^{1+2\alpha}}\,dx
+2|a-b|^2.
\label{eq:example-energy}
\end{align}
Accordingly, we introduce the Hilbert space
\[
\mathcal H_\alpha(\mathbb T)
=
\bigl\{u=(v,a,b):\ v\in H^\alpha(0,1),\ a,b\in\mathbb R\bigr\},
\]
endowed with the norm
\[
\|u\|_{\mathcal H_\alpha(\mathbb T)}^2
=
\|u\|_{L^2(\mu)}^2+\mathcal E_\alpha(u),
\]
where
\[
\|u\|_{L^2(\mu)}^2=\int_0^1 |v(x)|^2\,dx+a^2+b^2.
\]

\subsection{The variational problem}

Let $\lambda>0$, let $f\in L^2(0,1)$, and let $F,G\in\mathbb R$. We consider the functional
\begin{equation}\label{eq:example-functional}
\mathcal J(u)
=
\frac12\mathcal E_\alpha(u)
+\frac{\lambda}{2}\|u\|_{L^2(\mu)}^2
-
\left(
\int_0^1 f(x)v(x)\,dx+Fa+Gb
\right),
\end{equation}
for $u=(v,a,b)\in \mathcal H_\alpha(\mathbb T)$.

This is the nonlocal variational problem associated with the hybrid time scale $\mathbb T$:
\begin{equation}\label{eq:example-minimization}
\min_{u\in \mathcal H_\alpha(\mathbb T)} \mathcal J(u).
\end{equation}

\begin{proposition}\label{prop:example-existence}
The minimization problem \eqref{eq:example-minimization} admits a unique minimizer in $\mathcal H_\alpha(\mathbb T)$.
\end{proposition}

\begin{proof}
The functional $\mathcal J$ is the sum of the nonnegative quadratic form $\frac12\mathcal E_\alpha(u)$, the strictly positive quadratic term $\frac{\lambda}{2}\|u\|_{L^2(\mu)}^2$, and the continuous linear functional
\[
u=(v,a,b)\mapsto \int_0^1 f(x)v(x)\,dx+Fa+Gb.
\]
Hence $\mathcal J$ is strictly convex on $\mathcal H_\alpha(\mathbb T)$.

Moreover, by Cauchy--Schwarz,
\[
\left|\int_0^1 f(x)v(x)\,dx+Fa+Gb\right|
\le
\|f\|_{L^2(0,1)}\|v\|_{L^2(0,1)}+|F||a|+|G||b|
\le
C\|u\|_{L^2(\mu)}
\]
for some constant $C>0$. Therefore
\[
\mathcal J(u)
\ge
\frac12\mathcal E_\alpha(u)+\frac{\lambda}{2}\|u\|_{L^2(\mu)}^2-C\|u\|_{L^2(\mu)},
\]
so $\mathcal J$ is coercive on $\mathcal H_\alpha(\mathbb T)$. Since its quadratic part is convex and continuous and its linear part is weakly continuous, $\mathcal J$ is weakly lower semicontinuous. The direct method of the calculus of variations therefore yields the existence of a minimizer, and strict convexity implies uniqueness.
\end{proof}

\subsection{Weak formulation and continuous-discrete decomposition}

For
\[
u=(v,a,b)\in\mathcal H_\alpha(\mathbb T),
\qquad
\varphi=(\psi,\eta,\zeta)\in\mathcal H_\alpha(\mathbb T),
\]
we define the bilinear form
\begin{align}
\mathcal B_\alpha(u,\varphi)
&=
\int_0^1\int_0^1
\frac{(v(x)-v(y))(\psi(x)-\psi(y))}{|x-y|^{1+2\alpha}}\,dx\,dy
\notag\\
&\quad
+2\int_0^1
\frac{(a-v(x))(\eta-\psi(x))}{|2-x|^{1+2\alpha}}\,dx
\notag\\
&\quad
+2\int_0^1
\frac{(b-v(x))(\zeta-\psi(x))}{|3-x|^{1+2\alpha}}\,dx
+2(a-b)(\eta-\zeta).
\label{eq:example-bilinear}
\end{align}

\begin{proposition}\label{prop:example-weak}
A function $u=(v,a,b)\in\mathcal H_\alpha(\mathbb T)$ is the unique minimizer of $\mathcal J$ if and only if
\begin{equation}\label{eq:example-weak}
\mathcal B_\alpha(u,\varphi)
+\lambda\left(\int_0^1 v(x)\psi(x)\,dx+a\eta+b\zeta\right)
=
\int_0^1 f(x)\psi(x)\,dx+F\eta+G\zeta
\end{equation}
for every $\varphi=(\psi,\eta,\zeta)\in\mathcal H_\alpha(\mathbb T)$.
\end{proposition}

\begin{proof}
Let $u=(v,a,b)\in\mathcal H_\alpha(\mathbb T)$ and $\varphi=(\psi,\eta,\zeta)\in\mathcal H_\alpha(\mathbb T)$. For $t\in\mathbb R$,
\[
\mathcal J(u+t\varphi)
=
\frac12\mathcal E_\alpha(u+t\varphi)
+\frac{\lambda}{2}\|u+t\varphi\|_{L^2(\mu)}^2
-
\left(
\int_0^1 f(v+t\psi)\,dx+F(a+t\eta)+G(b+t\zeta)
\right).
\]
Differentiating with respect to $t$ at $t=0$, we obtain
\[
\frac{d}{dt}\mathcal J(u+t\varphi)\Big|_{t=0}
=
\mathcal B_\alpha(u,\varphi)
+\lambda\left(\int_0^1 v\psi\,dx+a\eta+b\zeta\right)
-\left(\int_0^1 f\psi\,dx+F\eta+G\zeta\right).
\]
Hence $u$ is a critical point of $\mathcal J$ if and only if \eqref{eq:example-weak} holds for every $\varphi$. By Proposition~\ref{prop:example-existence}, the critical point is unique and coincides with the unique minimizer.
\end{proof}

Equation \eqref{eq:example-weak} is the weak formulation of a hybrid nonlocal transmission problem on the time scale $\mathbb T$.

\begin{proposition}\label{prop:example-splitting}
Let $u=(v,a,b)\in\mathcal H_\alpha(\mathbb T)$. Then $u$ satisfies \eqref{eq:example-weak} for every $\varphi=(\psi,\eta,\zeta)\in\mathcal H_\alpha(\mathbb T)$ if and only if the following three conditions hold:

\medskip
\noindent
\textup{(i)} for every $\psi\in H^\alpha(0,1)$,
\begin{align}
&\int_0^1\int_0^1
\frac{(v(x)-v(y))(\psi(x)-\psi(y))}{|x-y|^{1+2\alpha}}\,dx\,dy
\notag\\
&\qquad
-2\int_0^1\frac{(a-v(x))\psi(x)}{|2-x|^{1+2\alpha}}\,dx
-2\int_0^1\frac{(b-v(x))\psi(x)}{|3-x|^{1+2\alpha}}\,dx
\notag\\
&\qquad
+\lambda\int_0^1 v(x)\psi(x)\,dx
=
\int_0^1 f(x)\psi(x)\,dx;
\label{eq:example-continuous}
\end{align}

\medskip
\noindent
\textup{(ii)}
\begin{equation}\label{eq:example-discrete-a}
2\int_0^1\frac{a-v(x)}{|2-x|^{1+2\alpha}}\,dx
+2(a-b)+\lambda a
=
F;
\end{equation}

\medskip
\noindent
\textup{(iii)}
\begin{equation}\label{eq:example-discrete-b}
2\int_0^1\frac{b-v(x)}{|3-x|^{1+2\alpha}}\,dx
+2(b-a)+\lambda b
=
G.
\end{equation}
\end{proposition}

\begin{proof}
Assume first that $u$ satisfies \eqref{eq:example-weak} for every $(\psi,\eta,\zeta)\in\mathcal H_\alpha(\mathbb T)$. Choosing $(\psi,\eta,\zeta)=(\psi,0,0)$ gives \eqref{eq:example-continuous}. Choosing $(\psi,\eta,\zeta)=(0,1,0)$ gives \eqref{eq:example-discrete-a}, and choosing $(\psi,\eta,\zeta)=(0,0,1)$ gives \eqref{eq:example-discrete-b}.

Conversely, assume that \eqref{eq:example-continuous}, \eqref{eq:example-discrete-a}, and \eqref{eq:example-discrete-b} hold. Let $(\psi,\eta,\zeta)\in\mathcal H_\alpha(\mathbb T)$ be arbitrary. Multiplying \eqref{eq:example-discrete-a} by $\eta$, multiplying \eqref{eq:example-discrete-b} by $\zeta$, and adding the resulting identities to \eqref{eq:example-continuous}, we recover exactly \eqref{eq:example-weak}. This proves the equivalence.
\end{proof}

The previous proposition shows that the hybrid weak problem consists of a nonlocal equation on the continuous phase, coupled with two scalar nonlocal balance relations on the discrete phase. The interface terms transmit information in both directions between the interval $[0,1]$ and the discrete nodes $2$ and $3$.

\subsection{A reduced Galerkin approximation}

The full variational problem above can be approximated on finite-dimensional subspaces of $\mathcal H_\alpha(\mathbb T)$. For instance, one may restrict the continuous component to affine functions,
\[
v(x)=mx+c,
\]
and consider the reduced space
\[
\mathcal X_{\mathrm{red}}
=
\{(mx+c,a,b):\ m,c,a,b\in\mathbb R\}.
\]
The restriction of $\mathcal J$ to $\mathcal X_{\mathrm{red}}$ then yields a finite-dimensional quadratic minimization problem, which may be viewed as a Galerkin approximation of the full hybrid transmission problem. This reduction is useful for explicit computations, but it should be regarded as a derived approximation of the full problem rather than as a substitute for it.
\section{Conclusion and Perspectives}

We have studied a hybrid nonlocal transmission problem on a continuous-discrete domain, formulated through a quadratic Gagliardo-type energy. The analysis shows that the interface contribution plays a central role: it yields a coercive coupling between the continuous and discrete phases and provides the key estimate underlying the whole variational theory. In particular, this leads to a Poincar\'e-type inequality adapted to the hybrid setting, as well as to the existence and uniqueness of a weak solution for the associated variational problem.

The weak formulation obtained in this way naturally splits into a nonlocal equation on the continuous component and discrete balance relations on the nodal part. This confirms that the model should be understood as a transmission problem rather than as a mere juxtaposition of continuous and discrete energies. The example discussed in the last section further illustrates that the hybrid framework remains genuinely variational and gives rise to explicit coupled systems.

A natural continuation of this work would be to consider more general finite discrete sets, to study reduced Galerkin-type approximations of the hybrid problem, and to extend the analysis to nonlinear energies with similar transmission structure. These directions remain close to the present setting and seem realistic within the same general approach.

\section*{Declarations}
\subsection*{competing interests} The authors declare that they have no known competing financial interests or personal relationships that could have appeared to influence the work reported in this paper.

\subsection*{Funding sources}
This research did not receive any specific grant from funding agencies in the public, commercial, or not-for-profit sectors.

\subsection*{Declaration of Generative AI and AI-assisted technologies in the writing process}
During the preparation of this work, the authors used ChatGPT to improve the language, readability, and presentation of the manuscript. The authors edited the content and take full responsibility for the final version of the manuscript.

\begin{thebibliography}{99}

\bibitem{AgarwalOteroPereraVivero2006}
R.~P. Agarwal, V.~Otero-Espinar, K.~Perera and D.~R. Vivero,
Basic properties of Sobolev's spaces on time scales,
\emph{Adv. Difference Equ.} \textbf{2006} (2006), Article ID 38121.

\bibitem{NonlocalWeakSolutions}
F.~Andreu, J.~M. Maz\'on, J.~D. Rossi, and J.~Toledo,
Existence and uniqueness of weak solutions for nonlocal parabolic problems,
\emph{J. Math. Anal. Appl.} \textbf{481} (2020), no.~1, Article 123437.

\bibitem{BenkhettouHassaniTorres2016}
N.~Benkhettou, S.~Hassani and D.~F.~M. Torres,
Fractional Sobolev's spaces on time scales via conformable fractional calculus and their application to a fractional differential equation,
\emph{Adv. Math. Phys.} \textbf{2016} (2016), Article ID 9636491.

\bibitem{BohnerGuseinov2006}
M.~Bohner and G.~Sh. Guseinov,
Multiple Lebesgue integration on time scales,
\emph{Adv. Difference Equ.} \textbf{2006} (2006), Article ID 26391.

\bibitem{BohnerPeterson2001}
M.~Bohner and A.~Peterson,
\emph{Dynamic Equations on Time Scales: An Introduction with Applications},
Birkh\"auser, Boston, 2001.

\bibitem{BohnerPeterson2003}
M.~Bohner and A.~Peterson,
\emph{Advances in Dynamic Equations on Time Scales},
Birkh\"auser, Boston, 2003.

\bibitem{CabadaVivero2006}
A.~Cabada and D.~R. Vivero,
Expression of the Lebesgue $\Delta$-integral on time scales as a usual Lebesgue integral. Application to the calculus of $\Delta$-antiderivatives,
\emph{Math. Comput. Modelling} \textbf{43} (2006), 194--207.

\bibitem{DiNezza2012}
E.~Di Nezza, G.~Palatucci and E.~Valdinoci,
Hitchhiker's guide to the fractional Sobolev spaces,
\emph{Bull. Sci. Math.} \textbf{136} (2012), 521--573.

\bibitem{Gagliardo1957}
E.~Gagliardo,
Caratterizzazioni delle tracce sulla frontiera relative ad alcune classi di funzioni in $n$ variabili,
\emph{Rend. Sem. Mat. Univ. Padova} \textbf{27} (1957), 284--305.

\bibitem{Hilger1990}
S.~Hilger,
Analysis on measure chains---a unified approach to continuous and discrete calculus,
\emph{Results Math.} \textbf{18} (1990), 18--56.

\bibitem{HuLi2022Left}
X.~Hu and Y.~Li,
Fractional Sobolev space on time scales and its application to a fractional boundary value problem,
\emph{Complexity} \textbf{2022} (2022), Article ID 7149356.

\bibitem{HuLi2022Right}
X.~Hu and Y.~Li,
Right fractional Sobolev space via Riemann--Liouville derivatives on time scales,
\emph{Fractal Fract.} \textbf{6} (2022), 121.

\bibitem{MRS2016}
G.~Molica Bisci, V.~D. R\u{a}dulescu, and R.~Servadei,
\emph{Variational Methods for Nonlocal Fractional Problems},
Encyclopedia of Mathematics and its Applications, Vol.~162,
Cambridge University Press, Cambridge, 2016.

\bibitem{SkrzypekSzymanska2019}
E.~Skrzypek and K.~Szyma'nska-D\k{e}bowska,
On the Lebesgue and Sobolev spaces on a time-scale,
\emph{Opuscula Math.} \textbf{39} (2019), 741--760.

\bibitem{Slobodeckij1958}
L.~N. Slobodeckij,
Generalized Sobolev spaces and their applications to boundary value problems for partial differential equations,
\emph{Leningrad. Gos. Ped. Inst. U\v cen. Zap.} \textbf{197} (1958), 54--112.

\bibitem{SteinbachWendland2000}
O.~Steinbach, W.~L. Wendland,
Variational formulations of transmission problems via FEM, BEM and DtN methods,
\emph{Comput. Struct.} \textbf{77} (2000), no.~1, 11--19.

\end{thebibliography}

\end{document}
